\documentclass[11pt]{article}
\usepackage[margin=1in]{geometry}
\usepackage[T1]{fontenc}
\usepackage{lmodern,amsmath,amssymb}
\usepackage[colorlinks=true,urlcolor=blue]{hyperref}
\title{Economic freedom and firm entry--exit decisions}
\author{Gabriel Saco\\Universidad del Pac\'ifico}
\date{Research status: 11 October 2026}
\begin{document}\maketitle
\noindent\textbf{Topic-stage progress document. This is not the final paper.}
The current four-page proposal is in \texttt{proposal/proposal.pdf}. The final
paper will require an 8--20-page body, complete proofs, simulations, and appendices.
\section*{Question and current model}
Can costly rebuilding make output losses after institutional deterioration larger
and more persistent than gains after improvement? The firm optimizes variable
input and activity under known temporary changes in usable property. The baseline
allows free restarting; the extension imposes a sunk restart cost.
\section*{Current result and its status}
The proposal states all maintained assumptions and the sufficient conditions
$\pi_0>0$, $\pi_-+\beta V<0$, and $\kappa>\pi_++\beta V$, where
$V=\pi_0/(1-\beta)$. Under these conditions the adverse incumbent closes and never
restarts, while the unchanged and favorable incumbents remain active. The current
proof is an outline, supported by numerical and symbolic checks. Lean
formalization and the handwritten appendix remain future work.
\section*{Remaining work}
I will complete the switching-path proof and handwritten derivations, study
suspension, entry, and heterogeneity, map the parameter region, and formalize
each numbered result using AppliedModelingLib. I will also clarify the scope of
aggregation and extend the literature review. Cheap suspension may eliminate the persistent loss.
\section*{Final-paper structure}
Introduction; related literature; numbered model assumptions; propositions and
proofs; reproducible simulations; limitations; conclusion; proof, Lean, and AI-log
appendices. \par The bibliography file records the current project sources.
\end{document}
